\subsection{SUM OF A FAMILY OF SETS}

PROPOSITION 9. Let \((\mathbf{X}_t)_{t \in I}\) be a family of sets. Then there exists a set \(\mathbf{X}\) with the following property: \(\mathbf{X}\) is the union of a family \((\mathbf{X}_t')_{t \in I}\) of mutually disjoint sets such that for each \(t \in I\) there exists a one-to-one mapping of \(\mathbf{X}_t\) onto \(\mathbf{X}_t'\) .

Let \(A = \bigcup_{i \in I} X_i\) . If \(i \in I\) , the mapping \(x \to (x, i)\) ( \(x \in X_i\) ) is a one-to-one mapping of \(X_i\) onto a subset \(X_i'\) of \(A \times I\) . Moreover, the image of \(X_i'\) under the second coordinate function on \(A \times I\) is contained in the set \(\{i\}\) ; it follows that \(X_i' \cap X_x' = \emptyset\) whenever \(i \neq x\) . We may then take \(X = \bigcup_{i \in I} X_i'\) .

DEFINITION 8. Let \((\mathbf{X}_{\iota})_{\iota \in \mathbf{I}}\) be a family of sets. The sum of this family is the union of the family of sets \((\mathbf{X}_{\iota} \times \{\iota\})_{\iota \in \mathbf{I}}\) .

PROPOSITION 10. Let \((\mathbf{X}_t)_{t\in I}\) be a family of mutually disjoint sets. Let A be its union and S its sum. Then there exists a bijective mapping of A onto S.

For each \(t \in I\) , let \(f_{t}\) be a bijection of \(X_{t}\) onto \(X_{t} \times \{t\}\) . By virtue of Proposition 8, there exists a mapping f of A into S which extends all the mappings \(f_{t}\) . It is immediately verified that f is a bijective mapping of A onto S.

By abuse of language, a set E is said to be the sum of a family of sets \((\mathbf{X}_{t})_{t\in\mathbf{I}}\) if there exists a bijection of E onto the sum of the family, as defined in Definition 8.

Note that if \((\mathbf{X}_{t})_{i\in\mathbf{I}}\) is an arbitrary family of sets, the argument of Proposition 10 shows that there exists a mapping of the sum S onto the union A.

